\chapter{The element identity story}
\label{ch:identity}

\begin{aside}
Two element trees that look the same may or may not be the same.
The renderer needs to know which: identical trees mean ``skip
the work''; different trees mean ``diff and emit''. The story
of how \maya{} establishes identity --- and why it matters ---
is more subtle than ``the trees are equal''.
\end{aside}

\section{Three notions of identity}

For an element tree, identity could mean:

\begin{itemize}[leftmargin=*]
  \item \textbf{Pointer identity.} Same address.
  \item \textbf{Structural identity.} Same shape, same contents.
  \item \textbf{Semantic identity.} Same shape, same contents,
    same focus / cursor / scroll positions.
\end{itemize}

Different renderer stages care about different ones. The cache
fingerprint cares about (2). The hit-tester cares about (1).
The animation system cares about (3).

\section{Pointer identity in the DSL}

The DSL builds a fresh tree every frame. The root pointer
changes; the children pointers change. Pointer identity is
useless for cross-frame reasoning.

If you need to keep state across frames keyed by ``the same
button'', use a stable ID:

\begin{lstlisting}
struct Button {
    int id;            // stable across frames
    std::string label;
};
\end{lstlisting}

The ID lives in your model. The element produced by rendering
the button gets the ID as a property; the runtime uses it for
focus, hit-test, animation state.

\section{Structural identity: the cache fingerprint}

The cache fingerprint (Chapter~\ref{ch:pipeline}) is structural
identity. Two trees with identical contents produce identical
fingerprints; two different trees produce different ones (with
high probability).

The fingerprint is computed by walking the tree and mixing each
leaf's content into a hash:

\begin{lstlisting}
uint64_t fingerprint(const Element& e) {
    XXH64State state;
    walk(e, [&](const TextElement& t) {
        XXH64_update(&state, t.text.data(), t.text.size());
        XXH64_update(&state, &t.style, sizeof(t.style));
    });
    walk(e, [&](const BoxElement& b) {
        XXH64_update(&state, &b.layout, sizeof(b.layout));
        // children handled by recursion
    });
    return XXH64_digest(&state);
}
\end{lstlisting}

The walk is depth-first, left-to-right. Two trees that differ
only in child order produce different fingerprints.

\subsection*{Collision probability}

A 64-bit hash collides at roughly the square-root rate: two
different trees out of $2^{32}$ are expected to share a
fingerprint. For a \maya{} app that produces at most a few
thousand distinct trees, the probability of a false ``no-change''
verdict in any given session is negligible.

If collision matters (e.g. you ship a long-running server with
millions of users), use a 128-bit hash. \maya{}'s fingerprint
is parameterised; switching is a one-line change.

\section{Semantic identity: keys and reconciliation}

When trees change, the renderer wants to preserve cursor
positions, scroll offsets, animation states. Naively, ``the
third child is now in position four'' would lose this state.

The fix: attach \emph{keys} to children. A key is a stable
identifier on an element that the reconciler uses to match
across frames.

\begin{lstlisting}
v(
    keyed(\"item-42\", render_item(items[42])),
    keyed(\"item-43\", render_item(items[43])),
    keyed(\"item-44\", render_item(items[44]))
);
\end{lstlisting}

The reconciler matches by key, not position. Inserting a new
item at the top does not shift the existing items' state.

\subsection*{When keys matter most}

\begin{itemize}[leftmargin=*]
  \item Lists where items can reorder.
  \item Lists with per-item animation state (slide-in,
    fade-out).
  \item Forms where input focus must follow a moving row.
\end{itemize}

For static lists, position-based reconciliation works fine.

\section{Identity in the layout phase}

Layout assigns each node a final rect. The rects are derived
from the tree, not from any cross-frame identity. Two frames
with the same tree produce the same rects.

This is intentional: layout is a pure function. There is no
``layout cache'' keyed by node identity. If you want the layout
to be different from frame to frame, the tree must differ.

\section{Identity in the canvas}

The canvas is a flat grid of cells. Cells have no identity
across frames; each frame produces a fresh canvas, and the diff
compares them cell-by-cell.

This means: cell-level animations (a single character pulsing)
are tracked at the cell level, not the element level. The
canvas does not know about elements.

\section{Identity in the focus system}

Focus is keyed by ID, not by element pointer. A focused button
keeps focus across re-renders because its ID is stable. If you
generate fresh IDs each frame, focus moves randomly --- the
classic ``why is focus jumping?'' bug.

\section{Identity in the signal system}

Signals are objects with their own identity. Subscribing to a
signal is keyed by the signal's address. A signal whose address
changes across frames cannot be subscribed to reliably; keep
signals in long-lived containers (the model, a long-lived
struct).

\section{When identity breaks down}

A few patterns cause identity to misbehave:

\subsection*{Anonymous lambdas as keys}

\begin{lstlisting}
auto button = [&] { return text(\"click\"); };
// Each frame, button is a different lambda --- different type, different identity.
\end{lstlisting}

Capture the lambda's expression, not the lambda itself. Or
better, write a function.

\subsection*{Index-based keys for reorderable lists}

\begin{lstlisting}
for (int i = 0; i < items.size(); ++i) {
    rows.push_back(keyed(std::to_string(i), render(items[i])));
}
\end{lstlisting}

When an item is inserted at the front, every key shifts. Use the
item's own stable ID, not the position.

\subsection*{Mutable state in element fields}

If an element carries mutable state via a captured reference,
the fingerprint may not detect changes (the data the reference
points to is not part of the element's bytes). Either include
the data by value in the element or mark the element as
``do not cache''.

\section{Pitfalls}

\begin{warning}
\textbf{Forgetting keys in reorderable lists.} Symptoms:
focus jumps to a different row, animation restarts, scroll
position drifts. Fix: add keys.
\end{warning}

\begin{warning}
\textbf{Keys that are not unique.} Two children with the same
key under one parent break the reconciler. Use IDs from
your model.
\end{warning}

\begin{warning}
\textbf{Cache fingerprint false negatives.} If your tree
references external mutable state, the fingerprint may not
change when the rendered output should. Make the data part
of the model, not external.
\end{warning}

\begin{warning}
\textbf{Cache fingerprint false positives.} If two visually
distinct trees somehow produce the same fingerprint
(collision), the renderer skips a real update. Vanishingly
rare with 64-bit hashes for typical app sizes; switch to
128-bit if you ship long-running services.
\end{warning}

\section*{Workshop}

\begin{enumerate}[leftmargin=*]
  \item Build a list of three items. Add a keyed item at the
    top. Confirm the bottom two items keep their state (focus,
    scroll).
  \item Replace the keys with the position index. Add another
    item at the top. Observe the state drift.
  \item Compute the fingerprint of a tree by hand (small one).
    Confirm two different trees yield different fingerprints.
\end{enumerate}

\begin{recap}
Identity in \maya{} has three notions: pointer (per frame),
structural (cross-frame, by content hash), and semantic
(cross-frame, by key). The renderer uses each appropriately.
Keys preserve focus, scroll, and animation state across
re-renders. The fingerprint enables ``skip everything'' on no
change. Mistakes in identity usually surface as visible state
loss.
\end{recap}
